\section{Conclusion}
\label{sec:conclusion}

S(M,T) assembles well-understood components (Product-of-Experts gates, additive utility, Boltzmann cost penalty, Robbins-Monro learning) into a scoring framework for routing across heterogeneous AI endpoints. The techniques are standard. The finding is not.

Fourteen attempts to compute hand-designed scoring functions from 2.76M benchmark records all failed. The reasons are structural: metrics use incomparable scales, benchmark performance does not transfer to deployment, and complex properties collapse when reduced to single numbers. This is the measurement gap.

Replacing hand-designed functions with 16 learned bilinear forms achieves 83.63\% accuracy on RouterBench and 94.35\% of best-model quality at half the cost on RouteLLM. The model generalizes across both benchmarks without benchmark-specific tuning. A 4.6$\times$ scaling experiment confirmed that data quality, not model capacity, is the binding constraint.

The measurement gap separates interpretable from learned routing, but it also points forward. If benchmarks cannot ground hand-designed routing scores, the field needs benchmarks designed for routing, not repurposed from model comparison. Close the measurement gap, and you do not just get better routing. You get a foundation for evaluating AI systems under the conditions where they actually operate.

Code: \texttt{github.com/pranavlakherwal/smt-router}

\paragraph{AI Disclosure.} AI tools from Anthropic assisted with writing and code development. All research direction, experimental design, and intellectual contributions are the author's own.
